\documentclass[10pt]{article}
\usepackage[ngerman]{babel}
\usepackage[utf8]{inputenc}
\usepackage[T1]{fontenc}
\usepackage{amsmath}
\usepackage{amsfonts}
\usepackage{amssymb}
\usepackage[version=4]{mhchem}
\usepackage{stmaryrd}
\usepackage{graphicx}
\usepackage[export]{adjustbox}
\graphicspath{ {./images/} }
\usepackage{bbold}

\begin{document}
Abiturprüfung an den allgemein bildenden Gymnasien

\section*{Prüfungsfach: Mathematik}
Beispielaufgabe 2024\\
Teil A\\
Blatt 1 von 4

\section*{Pflichtaufgaben}
Bearbeiten Sie alle Aufgaben P1 bis P4.

P1 Gegeben ist die in IR definierte Funktion $f$ mit $f(x)=\sin x$. Die Abbildung zeigt den Graphen $G_{f}$ von $f$ sowie die Tangenten an $G_{f}$ in den dargestellten Schnittpunkten mit der x-Achse.\\
\includegraphics[max width=\textwidth, center]{2025_04_15_6049cd5257d7821275b7g-01}\\
a Zeigen Sie, dass diejenige der beiden Tangenten, die durch den Koordinatenursprung verläuft, die Steigung 1 hat.\\
b Berechnen Sie den Inhalt des Flächenstücks, das von $G_{f}$ und den beiden Tangenten eingeschlossen wird.

P2 Die Verbreitung eines Computervirus lässt sich modellhaft mithilfe der in IR definierten Funktion $f$ mit $f(t)=2 \cdot t \cdot e^{-\frac{1}{100} t}$ beschreiben. Dabei ist $t$ die Zeit in Tagen, die seit der ersten Infizierung eines Computers mit dem Virus vergangen ist, und $f(t)$ die Rate der Infizierungen zum Zeitpunkt t in der Einheit „Eintausend Computer pro Tag".\\
a Zeigen Sie, dass $2 \cdot\left(1-\frac{1}{100} \mathrm{t}\right) \cdot \mathrm{e}^{-\frac{1}{100} \mathrm{t}}$ ein Term der ersten Ableitungsfunktion von $f$ ist.\\
b Geben Sie den Zeitpunkt an, zu dem die Rate der Infizierungen am größten ist.\\
c Betrachtet wird der Zeitraum der zweiten Woche nach der ersten Infizierung eines Computers mit dem Virus. Geben Sie einen Term an, mit dem die Anzahl der Computer berechnet werden kann, die in diesem Zeitraum infiziert werden.

P3 Gegeben sind der Punkt $P(-1|7| 2)$ und die Ebene $E: x_{1}+3 x_{2}=0$.\\
a Zeigen Sie, dass P nicht in E liegt.\\
b Bestimmen Sie die Koordinaten des Punkts, der entsteht, wenn $P$ an E gespiegelt wird.

P4 Die Vierfeldertafel gehört zu einem Zufallsexperiment mit Ereignissen $A$ und $B$. Für die Wahrscheinlichkeit $p$ gilt $p \neq 0$.

\begin{center}
\begin{tabular}{c|c|c|c}
 & $B$ & $\bar{B}$ &  \\
\hline
$A$ & $p$ &  & $3 p$ \\
\hline
$\bar{A}$ &  &  & $1-3 p$ \\
\hline
 & $4 p$ &  &  \\
\hline
\end{tabular}
\end{center}

a Vervollständigen Sie die Vierfeldertafel. Zeigen Sie, dass p nicht den Wert $\frac{1}{5}$ haben kann.\\
b Für einen bestimmten Wert von $p$ sind $A$ und $B$ stochastisch unabhängig. Ermitteln Sie diesen Wert von p.

\section*{Wahlaufgaben}
Bearbeiten Sie zwei der Aufgaben W1 bis W6.

W1 Gegeben sind die in IR definierten Funktionen $f$ und $g$. Der Graph von $f$ ist symmetrisch bezüglich der y-Achse, der Graph von g ist symmetrisch bezüglich des Koordinatenursprungs. Beide Graphen haben einen Hochpunkt im Punkt (2|1).\\
a Geben Sie für die Graphen von $f$ und $g$ jeweils die Koordinaten und die Art eines weiteren Extrempunkts an.\\
b Untersuchen Sie die in IR definierte Funktion $h$ mit $h(x)=f(x) \cdot(g(x))^{3}$ im Hinblick auf eine mögliche Symmetrie ihres Graphen.

W2 Die Abbildung zeigt den Graphen $G_{f}$ einer in $\mathbb{R}$ definierten Funktion f sowie den Graphen der ersten Ableitungsfunktion von f .\\
a Geben Sie die Steigung der Tangente an $G_{f}$ im Punkt (0|f(0)) an.\\
\includegraphics[max width=\textwidth, center]{2025_04_15_6049cd5257d7821275b7g-02}

Abiturprüfung an den allgemein bildenden Gymnasien

\section*{Prüfungsfach: Mathematik}
Beispielaufgabe 2024\\
Teil A\\
Blatt 3 von 4\\
b Betrachtet wird die Schar der Funktionen $g_{c}$ mit $c \in \mathbb{R}^{+}$. Der Graph von $g_{c}$ geht aus $G_{f}$ durch Streckung mit dem Faktor c in y -Richtung hervor. Die Tangente an den Graphen von $g_{c}$ im Punkt $\left(0 \mid g_{c}(0)\right)$ schneidet die x-Achse. Bestimmen Sie rechnerisch die x-Koordinate des Schnittpunkts.

W3 a Die Ebene E: $3 x_{1}+2 x_{2}+2 x_{3}=6$ enthält einen Punkt, dessen drei Koordinaten übereinstimmen. Bestimmen Sie diese Koordinaten.\\
b Begründen Sie, dass die folgende Aussage richtig ist:\\
Es gibt unendlich viele Ebenen, die keinen Punkt enthalten, dessen drei Koordinaten übereinstimmen.

W4 Gegeben sind die Punkte $A(0|0| 0), B(3|4| 1), C(1|7| 3)$ und $D(-2|3| 2)$.\\
a Weisen Sie nach, dass das Viereck ABCD ein Parallelogramm ist.\\
b Der Punkt T liegt auf der Strecke $\overline{A C}$. Das Dreieck $A B T$ hat bei $B$ einen rechten Winkel. Ermitteln Sie das Verhältnis der Länge der Strecke $\overline{\mathrm{AT}}$ zur Länge der Strecke $\overline{\mathrm{CT}}$.

W5 In einem Behälter befinden sich Kugeln, von denen jede dritte gelb ist.\\
a Aus dem Behälter wird zweimal nacheinander jeweils eine Kugel zufällig entnommen und wieder zurückgelegt. Berechnen Sie die Wahrscheinlichkeit dafür, dass beide Kugeln gelb sind.\\
b Im Behälter werden zwei gelbe Kugeln durch zwei blaue Kugeln ersetzt. Anschließend wird aus dem Behälter erneut zweimal nacheinander jeweils eine Kugel zufällig entnommen und wieder zurückgelegt. Die Wahrscheinlichkeit dafür, dass beide Kugeln gelb sind, beträgt nun $\frac{1}{16}$. Ermitteln Sie, wie viele gelbe Kugeln sich nach den beschriebenen Vorgängen im Behälter befinden.

\section*{Prüfungsfach: Mathematik}
W6 Für ein Spiel wird ein Glücksrad verwendet. Die Abbildung zeigt dieses Glücksrad schematisch. Die Wahrscheinlichkeit dafür, bei einmaligem Drehen die Zahl 2 zu erzielen, wird mit $p$ bezeichnet.

Bei dem Spiel bezahlt jeder Spieler zunächst einen Einsatz von 1 Euro. Anschließend dreht er das Glücksrad zweimal. Erzielt er dabei zwei Zahlen, deren Summe mindestens 5 ist, wird inm der Wert der Summe als Betrag in Euro ausgezahlt;\\
\includegraphics[max width=\textwidth]{2025_04_15_6049cd5257d7821275b7g-04} ansonsten erfolgt keine Auszahlung. Wird das Spiel wiederholt durchgeführt, so ist zu erwarten, dass sich auf lange Sicht die Einsätze der Spieler und die Auszahlungen ausgleichen.

Leiten Sie unter Verwendung der beschriebenen Spielregeln eine Gleichung her, mit der der Wert von p berechnet werden könnte; erläutern Sie dabei Ihr Vorgehen.

Abiturprüfung an den allgemein bildenden Gymnasien\\
Prüfungsfach: Mathematik\\
Aufgabe I 1\\
Beispielaufgabe 2024\\
Teil B: Analysis\\
Blatt 1 von 2

Die Abbildung 1 zeigt modellhaft den Längsschnitt einer dreiteiligen Brücke aus Holz für eine Spielzeugeisenbahn. Die Züge können sowohl über die Brücke fahren als auch darunter hindurch.\\
\includegraphics[max width=\textwidth, center]{2025_04_15_6049cd5257d7821275b7g-05}

Länge der Brücke\\
Abb. 1\\
Die obere Randlinie des Längsschnitts der Brücke kann mithilfe des Graphen der in IR definierten Funktion $f$ mit $f(x)=\frac{1}{20} x^{4}-\frac{2}{5} x^{2}+1$ beschrieben werden. Dabei werden die Endpunkte dieser Randlinie durch die beiden Tiefpunkte des Graphen von $f$ dargestellt. Im verwendeten Koordinatensystem beschreibt die x-Achse die Horizontale; eine Längeneinheit entspricht einem Dezimeter in der Realität.

1 a Zeigen Sie rechnerisch, dass die obere Randlinie achsensymmetrisch ist.\\
b Bestimmen Sie rechnerisch die Höhe und die Länge der Brücke.\\
(zur Kontrolle: Ein Tiefpunkt des Graphen von f hat die x-Koordinate 2.)\\
c Betrachtet wird derjenige Punkt der oberen Randlinie, der sich am Übergang vom mittleren zum rechten Bauteil befindet. Prüfen Sie, ob dieser Punkt auf halber Höhe zwischen dem höchsten Punkt der oberen Randlinie und deren rechtem Endpunkt liegt.\\
d Geben Sie die Bedeutung des Terms $\frac{f(2)-f(1)}{2-1}$ im Sachzusammenhang an und berechnen Sie seinen Wert.\\
e Berechnen Sie die Größe des größten Steigungswinkels der Brücke, der beim Überfahren zu überwinden ist.

Der parabelförmige Teil der unteren Randlinie des Längsschnitts der Brücke kann mithilfe des Graphen einer in $\mathbb{R}$ definierten Funktion $q$ mit $q(x)=0,8-a \cdot x^{2}$ mit $a \in \mathbb{R}^{+}$beschrieben werden.

Abiturprüfung an den allgemein bildenden Gymnasien\\
f In der Abbildung 1 ist die Länge einer der beiden Bodenflächen des mittleren Bauteils mit s bezeichnet. Bestimmen Sie alle Werte von a, die für diese Länge mindestens $0,1 \mathrm{dm}$ liefern.\\
g Begründen Sie im Sachzusammenhang, dass für die Beschreibung der unteren Randlinie beliebig große Werte von a nicht infrage kommen.\\
h Für die Brücke gilt $\mathbf{a}=1,25$. Die drei Bauteile der Brücke werden aus massivem Holz hergestellt; $1 \mathrm{dm}^{3}$ des Holzes hat eine Masse von 800 Gramm. Die Brücke ist $0,4 \mathrm{dm}$ breit. Ermitteln Sie die Masse des mittleren Bauteils.

2 Während der Planung der Brückenform kamen zur Beschreibung der oberen Randlinie für das linke Bauteil eine Funktion $g_{\ell}$ und für das rechte Bauteil eine Funktion $g_{r}$ infrage. Auch bei Verwendung dieser Funktionen wäre die obere Randlinie achsensymmetrisch gewesen. Beurteilen Sie jede der folgenden Aussagen:


\begin{equation*}
\text { I }-g_{\ell}(x)=g_{r}(-x) \text { für }-2 \leq x \leq-1 \quad \text { II } g_{\ell}(x-1)=g_{r}(-x+1) \text { für }-1 \leq x \leq 0 \tag{4BE}
\end{equation*}


3 Die Form und die Größe der Brücke werden verändert, indem im bisher verwendeten Modell die obere Randlinie des Längsschnitts mithilfe der in IR definierten Funktion $k$ mit $k(x)=\frac{3}{5} \cdot \cos \left(\frac{\pi}{3} x\right)+\frac{4}{5}$ beschrieben wird. Die Bauteile der veränderten Brücke lassen sich nach dem in der Abbildung 2 dargestellten Prinzip aus einem quaderförmigen Holzblock sägen. Der beim Sägen auftretende Materialverlust soll im Folgenden vernachlässigt werden.\\
\includegraphics[max width=\textwidth, center]{2025_04_15_6049cd5257d7821275b7g-06}

Abb. 2\\
a Der Graph von $k$ ist symmetrisch bezüglich jedes seiner Wendepunkte. Beschreiben Sie, wie diese Eigenschaft mit dem in der Abbildung 2 dargestellten Prinzip zusammenhängt.\\
b Ermitteln Sie mithilfe des Funktionsterms von k den Flächeninhalt der gesamten in der Abbildung 2 gezeigten rechteckigen Vorderseite des Holzblocks.

Abiturprüfung an den allgemein bildenden Gymnasien\\
Prüfungsfach: Mathematik\\
Aufgabe 12\\
Beispielaufgabe 2024\\
Teil B: Analysis\\
Blatt 1 von 3

Gegeben ist die in $\mathbb{R}$ definierte Funktion $f$ mit $f(x)=e^{x}$.\\
a Bestimmen Sie rechnerisch eine Gleichung der Tangente an den Graphen von f in dessen Schnittpunkt mit der y-Achse. Geben Sie eine Gleichung der Gerade an, die in diesem Punkt senkrecht zur Tangente steht, sowie den Inhalt der Fläche, die von dieser Gerade, der Tangente und der x-Achse eingeschlossen wird.\\
b Geben Sie alle Werte von $m \in \mathbb{R}$ an, für die die Gerade mit der Gleichung $y=m \cdot x+1$ mit dem Graphen von $f$ genau einen Punkt gemeinsam hat.\\
c Der Graph von f, die Koordinatenachsen und die Gerade mit der Gleichung $\mathrm{x}=\mathrm{a}$ mit a>0 begrenzen ein Flächenstück. Rotiert dieses Flächenstück um die $x$-Achse, so erzeugt es einen Körper mit dem Volumen $3 \pi$. Bestimmen Sie den Wert von a.

Gegeben sind außerdem die in $\mathbb{R} \backslash\{0\}$ definierte Funktion $g$ mit $g(x)=\frac{1}{x}$ sowie die Funktion $h$ mit $h(x)=\frac{1}{x-2}+4$, deren Definitionsmenge so groß wie möglich gewählt wurde.\\
d Beschreiben Sie, wie der Graph von h aus dem Graphen von g erzeugt werden kann. Geben Sie die Definitionsmenge von h sowie den Grenzwert von h für $x \rightarrow+\infty$ an.\\
e Die Abbildung in der Anlage soll den Graphen von h sowie die Gerade w mit der Gleichung $y=x$ darstellen. Ergänzen Sie die fehlenden Koordinatenachsen und skalieren Sie diese passend.\\
f Für jede Stammfunktion H von $h$ gilt für jede reelle Zahl $\mathrm{b}>5$ :\\
$\mathrm{H}(\mathrm{b})-\mathrm{H}(5)>4 \cdot(\mathrm{~b}-5)$. Beschreiben Sie die geometrische Bedeutung dieser Aussage und veranschaulichen Sie die Aussage grafisch.

Die Umkehrfunktion von h wird mit k bezeichnet.\\
g Alle gemeinsamen Punkte der Graphen von $h$ und $k$ liegen auf der Gerade w. Berechnen Sie die x-Koordinaten dieser gemeinsamen Punkte.\\
h Begründen Sie ohne zu rechnen, dass jede Tangente mit der Steigung -1 an den Graphen von $h$ auch Tangente an den Graphen von kist.

Prüfungsfach: Mathematik\\
Beispielaufgabe 2024

Aufgabe 12\\
Blatt 2 von 3

Betrachtet werden nun die in $\mathbb{R} \backslash\{0\}$ definierte Funktion u mit $u(x)=f(g(x))$ sowie die in $\mathbb{R}$ definierte Funktion $v(x)=g(f(x))$. Für die erste Ableitungsfunktion $u^{\prime}$ von $u$ gilt $u^{\prime}(x)=-\frac{1}{x^{2}} \cdot e^{\frac{1}{x}}$.\\
i Geben Sie für jede der Funktionen u und v einen Funktionsterm an, der zwar die Variable x, aber keine Funktionsbezeichnung wie "f" oder "g" enthält. Nennen Sie die Wertemengen von $u$ und $v$.\\
j Zeigen Sie, dass u genau eine Wendestelle hat, und bestimmen Sie diese Wendestelle.

Prüfungsfach: Mathematik (Beispielaufgabe 2024)

Zu- und Vorname: $\qquad$\\
Chiffre der Schule\\
Blatt 3 von 3\\
Chiffre des Schülers $\square$

Prüfungsfach: Mathematik (Beispielaufgabe 2024)

Abbildung zu Teilaufgabe e

\begin{center}
\begin{tabular}{|c|c|c|c|c|c|c|c|c|c|c|c|c|c|}
\hline
 &  &  &  &  &  &  &  &  &  &  & ,' &  &  \\
\hline
 &  &  &  &  &  &  &  &  &  & ,' &  &  &  \\
\hline
 &  &  &  &  &  &  &  &  & , ' &  &  &  &  \\
\hline
 &  &  &  & $\underline{ }$ &  &  &  & , ' &  &  &  &  &  \\
\hline
 &  &  &  & $\bigcirc$ &  &  & ,' &  &  &  &  &  &  \\
\hline
 &  &  &  &  &  &  &  &  &  &  &  &  &  \\
\hline
 &  &  &  &  & ,'' &  &  &  &  &  &  &  &  \\
\hline
 &  &  &  & \( ,^{\prime} \) &  &  &  &  &  &  &  &  &  \\
\hline
 &  &  & , &  &  &  &  &  &  &  &  &  &  \\
\hline
 &  & , ' &  &  &  &  &  &  &  &  &  &  &  \\
\hline
 & ,' &  &  &  &  &  &  &  &  &  &  &  &  \\
\hline
1,' &  &  &  &  &  &  &  &  &  &  &  &  &  \\
\hline
\end{tabular}
\end{center}

Abiturprüfung an den allgemein bildenden Gymnasien

Teil B: Analytische Geometrie\\
Blatt 1 von 2

In einem Koordinatensystem wird der abgebildete Körper ABCDEF mit $A(0|10| 1), B(10|20| 1)$, $C(0|20| 1), D(0|7| 0)$ und $F(0|20| 0)$ betrachtet. Die beiden Seitenflächen ACFD und BEFC stehen senkrecht zur $\mathrm{x}_{1} \mathrm{x}_{2}$-Ebene.\\
\includegraphics[max width=\textwidth, center]{2025_04_15_6049cd5257d7821275b7g-10}

1 a $A, B$ und $D$ liegen in der Ebene $H$. Bestimmen Sie eine Gleichung von H in Koordinatenform.\\
(zur Kontrolle: $x_{1}-x_{2}+3 x_{3}+7=0$ )\\
b Begründen Sie, dass die Gerade $\mathrm{i}: \overrightarrow{\mathrm{x}}=\left(\begin{array}{l}0 \\ 7 \\ 0\end{array}\right)+\lambda \cdot\left(\begin{array}{l}1 \\ 1 \\ 0\end{array}\right)$ mit $\lambda \in \mathbb{R}$ sowohl in der $x_{1} x_{2}$-Ebene als auch in der Ebene H liegt.

Der Punkt E liegt auf $i$, wobei der Abstand von E zu F ebenso groß ist wie der Abstand von D zu F.\\
c Ermitteln Sie die Koordinaten von E.\\
d Begründen Sie ohne zu rechnen, dass die Vierecke ACFD und BEFC den gleichen Flächeninhalt haben.

2 Der Körper ABCDEF stellt modellhaft ein Podest dar, das auf der Bühne eines Theaters steht, das Viereck ADEB die Vorderseite des Podests und der Punkt D deren untere linke Ecke. Die $x_{1} x_{2}$-Ebene beschreibt den horizontalen Boden der Bühne. Eine Längeneinheit im Koordinatensystem entspricht einem Meter in der Wirklichkeit.\\
a Zeigen Sie, dass die Deckfläche des Podests rechtwinklig ist, und berechnen Sie deren Flächeninhalt.\\
b Die Position eines Scheinwerfers kann im Modell durch den Punkt $S(5|-3| z)$ dargestellt werden. Vom Scheinwerfer ausgehendes Licht trifft an der unteren linken Ecke der Vorderseite des Podests unter einem Winkel der Größe $47^{\circ}$ auf den Boden auf. Ermitteln Sie die Höhe des Scheinwerfers über dem Boden der Bühne.\\
c Die Position eines zweiten Scheinwerfers lässt sich im Modell durch den Punkt $P(2|4| 8)$ beschreiben. Die Gerade mit der Gleichung\\
$\vec{x}=\left(\begin{array}{c}0 \\ 10 \\ 1\end{array}\right)+\mu \cdot\left(\begin{array}{l}1 \\ 1 \\ 0\end{array}\right)$ und $\mu \in \mathbb{R}$ schneidet die Ebene mit der Gleichung\\
$\left(\overrightarrow{\mathrm{x}}-\left(\begin{array}{l}2 \\ 4 \\ 8\end{array}\right)\right) \circ\left(\begin{array}{l}1 \\ 1 \\ 0\end{array}\right)=0$ im Punkt $\mathrm{Q}(-2|8| 1)$. Es gilt $|\overrightarrow{\mathrm{PQ}}|=9$.\\
Treffen Sie auf der Grundlage dieser Informationen eine Aussage über den Abstand des zweiten Scheinwerfers von der Vorderkante der Deckfläche des Podests. Begründen Sie Ihre Aussage ohne zu rechnen.

Abiturprüfung an den allgemein bildenden Gymnasien\\
Prüfungsfach: Mathematik\\
Aufgabe II 2\\
Beispielaufgabe 2024\\
Teil B: Analytische Geometrie\\
Blatt 1 von 1

Gegeben sind die Punkte $A(5|-5| 12), B(5|5| 12)$ und $C(-5|5| 12)$.\\
a Zeigen Sie, dass das Dreieck ABC gleichschenklig ist.\\
b Begründen Sie, dass A, B und C Eckpunkte eines Quadrats sein können, und geben Sie die Koordinaten des vierten Eckpunkts D dieses Quadrats an.

Im Folgenden wird die abgebildete Doppelpyramide betrachtet. Die beiden Teilpyramiden ABCDS und ABCDT sind gleich hoch. Der Punkt T liegt im Koordinatenursprung, der Punkt S ebenfalls auf der $x_{3}$-Achse.\\
Die Seitenfläche BCT liegt in einer Ebene E.\\
c Bestimmen Sie eine Gleichung von E in Koordinatenform.\\
\includegraphics[max width=\textwidth, center]{2025_04_15_6049cd5257d7821275b7g-12}\\
(zur Kontrolle: $12 x_{2}-5 x_{3}=0$ )\\
d Bestimmen Sie die Größe des Winkels, den die Seitenfläche BCT mit der Fläche ABCD einschließt.\\
$E$ gehört zur Schar der Ebenen $E_{k}: k x_{2}-5 x_{3}=5 k-60$ mit $k \in \mathbb{R}$.\\
e Alle Ebenen der Schar schneiden sich in einer Gerade. Weisen Sie nach, dass die Kante $\overline{B C}$ auf dieser Gerade liegt.\\
f Ermitteln Sie diejenigen Werte von $k$, für die $E_{k}$ mit der Seitenfläche ADS mindestens einen Punkt gemeinsam hat.\\
g Die Seitenfläche ADT liegt in der Ebene F. Geben Sie einen Normalenvektor von F an und begründen Sie Ihre Angabe, ohne die Koordinaten von $A$ und $D$ zu verwenden. Bestimmen Sie denjenigen Wert von $k$, für den $E_{k}$ senkrecht zu $F$ steht.\\
h Die Doppelpyramide wird so um die $x_{1}$-Achse gedreht, dass die bisher mit BCT bezeichnete Seitenfläche in der $x_{1} x_{2}$-Ebene liegt und der bisher mit $S$ bezeichnete Punkt eine positive $x_{2}$-Koordinate hat. Bestimmen Sie diese $x_{2}$-Koordinate und veranschaulichen Sie Ihr Vorgehen durch eine Skizze.

Abiturprüfung an den allgemein bildenden Gymnasien\\
Prüfungsfach: Mathematik\\
Aufgabe III 1\\
Beispielaufgabe 2024\\
Teil B: Stochastik\\
Blatt 1 von 2

Ein Unternehmen produziert Stahlkugeln für Kugellager. Erfahrungsgemäß sind $4 \%$ aller Kugeln fehlerhaft.

800 Kugeln werden zufällig ausgewählt. Die Anzahl der fehlerhaften Kugeln unter den ausgewählten kann durch eine binomialverteilte Zufallsgröße beschrieben werden.\\
a Bestimmen Sie die Wahrscheinlichkeit dafür, dass unter den ausgewählten Kugeln weniger als 30 fehlerhaft sind.\\
b Bestimmen Sie die Wahrscheinlichkeit dafür, dass die Anzahl der fehlerhaften Kugeln unter den ausgewählten höchstens um eine halbe Standardabweichung vom Erwartungswert dieser Anzahl abweicht.

Eine fehlerhafte Kugel hat entweder einen Formfehler oder einen Größenfehler. Die Wahrscheinlichkeit dafür, dass eine Kugel einen Formfehler hat, beträgt 3 \%. Alle Kugeln werden vor dem Verpacken geprüft. Dabei werden 95 \% der Kugeln mit Formfehler, 98 \% der Kugeln mit Größenfehler, aber auch 0,5 \% der Kugeln ohne Fehler aussortiert.\\
c Stellen Sie den Sachzusammenhang in einem beschrifteten Baumdiagramm dar.\\
d Berechnen Sie die Wahrscheinlichkeit dafür, dass eine aussortierte Kugel keinen Formfehler hat.

Aufgrund zunehmender Reklamationen wird vermutet, dass der Anteil der fehlerhaften Kugeln auf über $4 \%$ angestiegen ist. Um diese Vermutung zu prüfen, soll die Nullhypothese „Der Anteil der fehlerhaften Kugeln beträgt höchstens 4 \%." auf der Grundlage einer Stichprobe von 500 Kugeln getestet werden. Wenn das Ergebnis des Tests die Vermutung nicht entkräftet, soll die Produktion unterbrochen werden, um die Maschinen zu warten. Das Risiko, die Produktion irrtümlich zu unterbrechen, soll höchstens $3 \%$ betragen.\\
e Beschreiben Sie für diesen Test im Sachzusammenhang den Fehler zweiter Art. Geben Sie die Konsequenz an, die sich aus diesem Fehler für die Produktion ergeben würde.\\
f Für den beschriebenen Test wird der Ablehnungsbereich betrachtet. Eine der beiden Grenzen dieses Ablehnungsbereichs ist größer als 0 und kleiner als 500; diese Grenze wird mit $k$ bezeichnet. Zur Bestimmung des Werts von $k$ soll die binomial-

Prüfungsfach: Mathematik\\
Beispielaufgabe 2024\\
Teil B: Stochastik\\
Aufgabe III 1\\
Blatt 2 von 2\\
verteilte Zufallsgröße Y mit den Parametern $n=500$ und $p=0,04$ verwendet werden. Begründen Sie, dass keine der beiden Ungleichungen I und II den korrekten Wert von k liefert.\\
I $\mathrm{P}(\mathrm{Y} \leq \mathrm{k}) \leq 0,03$\\
II $\mathrm{P}(\mathrm{Y} \leq \mathrm{k}) \geq 0,97^{*}$\\
g Die Kugeln werden in Packungen verkauft. Ein Teil der verkauften Packungen wird zurückgegeben. Die Wahrscheinlichkeit dafür, dass eine verkaufte Packung zurückgegeben wird, beträgt $3 \%$. Dem Unternehmen entsteht pro Packung, die zurückgegeben wird, ein Verlust von 5,80 Euro; pro Packung, die nicht zurückgegeben wird, erzielt das Unternehmen einen Gewinn von 8,30 Euro. Bestimmen Sie die Wahrscheinlichkeit, mit der das Unternehmen bei einem Verkauf von 200 Packungen einen Gesamtgewinn von mindestens 1500 Euro erzielt.

Abiturprüfung an den allgemein bildenden Gymnasien

Prüfungsfach: Mathematik\\
Beispielaufgabe 2024\\
Teil B: Stochastik\\
Aufgabe III 2\\
Blatt 1 von 2

1 Ein Fahrradhändler hat festgestellt, dass es sich bei 40 \% aller von inm verkauften Fahrräder um Mountainbikes handelt. Es soll davon ausgegangen werden, dass in einer zufälligen Auswahl verkaufter Fahrräder die Anzahl der Mountainbikes binomialverteilt ist.\\
a Bestimmen Sie die Wahrscheinlichkeit dafür, dass sich in einer zufälligen Auswahl von 100 verkauften Fahrrädern genau 30 Mountainbikes befinden.\\
b Beschreiben Sie im Sachzusammenhang ein Zufallsexperiment, bei dem die Wahrscheinlichkeit eines Ereignisses mit dem Term 1-0,6 ${ }^{10}-10 \cdot 0,4 \cdot 0,6^{9}$ berechnet werden kann. Geben Sie dieses Ereignis an.\\
c Ermitteln Sie die Wahrscheinlichkeit dafür, dass in einer zufälligen Auswahl von 250 verkauften Fahrrädern die Anzahl der Mountainbikes um mindestens 10 \% größer ist als der Erwartungswert für diese Anzahl.\\
d $20 \%$ aller verkauften Fahrräder haben einen Rahmen aus Aluminium. $45 \%$ aller verkauften Fahrräder sind weder Mountainbikes noch haben sie einen Rahmen aus Aluminium. Bestimmen Sie den Anteil der Fahrräder mit einem Rahmen aus Aluminium unter den verkauften Mountainbikes.\\
e Der Händler hat berechnet, dass er im September des Jahres 2020 mit einer Wahrscheinlichkeit von mehr als 96 \% mehr als 800 Mountainbikes verkaufen wird. Ermitteln Sie, von welcher Anzahl verkaufter Fahrräder er bei seiner Berechnung mindestens ausgegangen ist.

Die Abbildung zeigt für einige Monate des Jahres 2019 jeweils den Anteil der Mountainbikes unter allen verkauften Fahrrädern.\\
\includegraphics[max width=\textwidth, center]{2025_04_15_6049cd5257d7821275b7g-15}\\
f Im April wurden 810 Mountainbikes verkauft. Bestimmen Sie für diesen Monat die Anzahl aller verkauften Fahrräder.\\
g Der Anteil der Mountainbikes lag im Mai und Juni insgesamt bei 46 \%; im Juli war er größer als im Mai und im August größer als im Juni. Entscheiden Sie, ob es dennoch möglich ist, dass der Anteil der Mountainbikes im Juli und August insgesamt kleiner war als insgesamt im Mai und Juni. Begründen Sie Ihre Entscheidung.

2 Zucker wird in unterschiedlich großen Packungen angeboten. Es soll davon ausgegangen werden, dass für jede Packungsgröße die tatsächliche Masse des Zuckers durch eine normalverteilte Zufallsgröße beschrieben werden kann.\\
a Für eine bestimmte Packungsgröße ist $\varphi(x)=\frac{1}{\sqrt{200 \pi}} \cdot \mathrm{e}^{-\frac{(x-500)^{2}}{200}}$ der Term der zugehörigen Dichtefunktion, wobei $x$ die Masse des Zuckers in Gramm ist. Geben Sie den Erwartungswert und die Standardabweichung für die Masse jeweils in Gramm an.\\
b Bei einer anderen Packungsgröße beträgt der Erwartungswert für die Masse des Zuckers 250 g , die Standardabweichung 5 g . Bestimmen Sie - auf 1 g genau - das kleinste Intervall, das mit einer Wahrscheinlichkeit von mindestens 98 \% die tatsächliche Masse des Zuckers enthält und symmetrisch bezüglich des Erwartungswerts ist.


\end{document}